\documentclass[conference]{IEEEtran}

% ── Packages ──────────────────────────────────────────────────────────────────
\usepackage[T1]{fontenc}
\usepackage[utf8]{inputenc}
\usepackage{amsmath,amssymb,amsfonts}
\usepackage{graphicx}
\usepackage{booktabs}
\usepackage{tabularx}
\usepackage{multirow}
\usepackage{multicol}
\usepackage{algorithm}
\usepackage{algpseudocode}
\usepackage{cite}
\usepackage{url}
\usepackage{hyperref}
\usepackage{xcolor}
\usepackage{array}
\usepackage{float}

% ── Fix IEEEtran algorithm float ──────────────────────────────────────────────
\makeatletter
\newif\if@restonecol
\makeatother

% ── Column type for tables ────────────────────────────────────────────────────
\newcolumntype{C}[1]{>{\centering\arraybackslash}p{#1}}
\newcolumntype{L}[1]{>{\raggedright\arraybackslash}p{#1}}

% ── Title block ───────────────────────────────────────────────────────────────
\title{Comparison of Apriori Algorithm for Frequent Itemset Mining with\\
ProbDF (2024) and Optimization Strategies for Improved Performance}

\author{
  \IEEEauthorblockN{Bilal Ahmad Sheikh}
  \IEEEauthorblockA{
    Department of Artificial Intelligence\\
    Ghulam Ishaq Khan Institute\\
    Topi, Pakistan\\
    2023162@giki.edu.pk
  }
  \and
  \IEEEauthorblockN{Muhammad Muzzammil Idrees}
  \IEEEauthorblockA{
    Department of Artificial Intelligence\\
    Ghulam Ishaq Khan Institute\\
    Topi, Pakistan\\
    2023481@giki.edu.pk
  }
  \and
  \IEEEauthorblockN{Raneem Sultan}
  \IEEEauthorblockA{
    Department of Artificial Intelligence\\
    Ghulam Ishaq Khan Institute\\
    Topi, Pakistan\\
    2023597@giki.edu.pk
  }
  \and
  \IEEEauthorblockN{~}
  \IEEEauthorblockA{~}
  \and
  \IEEEauthorblockN{Muhammad Uzair}
  \IEEEauthorblockA{
    Department of Artificial Intelligence\\
    Ghulam Ishaq Khan Institute\\
    Topi, Pakistan\\
    2023547@giki.edu.pk
  }
  \and
  \IEEEauthorblockN{~}
  \IEEEauthorblockA{~}
}

\begin{document}
\maketitle

% ─────────────────────────────────────────────────────────────────────────────
% AUTHOR CONTRIBUTIONS
% ─────────────────────────────────────────────────────────────────────────────
\begin{IEEEkeywords}
\textbf{Author Contributions:}
Bilal Ahmad Sheikh led implementation of Apriori Baseline, HashTree
optimization, and ProbDF, and produced Notebooks~1--4 (data loading,
benchmarking, comparison, visualization).
Muhammad Muzzammil Idrees implemented the bitset Tidset optimization,
conducted the exact precision/recall validation analysis, and co-authored
the literature review and algorithm description sections.
Raneem Sultan designed the experimental protocol, produced the comparison
figures and tractability map, and authored the experimental-setup and
discussion sections.
Muhammad Uzair derived the complexity-analysis tables, authored the
optimization-strategy justifications, conclusion, and reference list.
All four members contributed equally to viva preparation and are
prepared to discuss all aspects of the project.
\end{IEEEkeywords}

% ─────────────────────────────────────────────────────────────────────────────
% ABSTRACT
% ─────────────────────────────────────────────────────────────────────────────
\begin{abstract}
Frequent itemset mining (FIM) is a foundational task in data mining that
identifies co-occurring item sets in transactional databases, serving as
the backbone of association rule mining. Despite decades of research, the
classical Apriori algorithm remains computationally prohibitive on dense,
real-world datasets due to exponential candidate generation and repeated
database scans. This paper presents a rigorous empirical comparison of
Apriori against ProbDF (Sadeequllah et al., 2024), a state-of-the-art
approximate FIM algorithm that stops reading transaction data after
computing 2-itemsets and employs a Probabilistic Support Prediction Model
(PSPM) for higher-order itemsets. We additionally propose and evaluate two
optimization strategies for Apriori: a hash-tree based candidate lookup
(Park et al., 1995) and a bitset-accelerated vertical tidset format
(Zaki et al., 1997). Experiments are conducted on three FIMI benchmark
datasets---Chess, Connect, and Accidents---across five minimum support
thresholds (50\%--90\%). Results demonstrate that ProbDF achieves up to
11.63$\times$ speedup over Apriori Baseline on Chess with memory usage
under 130\,MB across all tested Chess thresholds, while Apriori variants
become intractable below 80\% support on dense datasets. The Connect
dataset proved entirely intractable for all Apriori variants, which itself
constitutes a key finding motivating approximate methods. Exact validation
at Chess@90\% shows ProbDF achieves 99.1\% precision and 91.3\% recall
against ground truth, with recall ranging from 85\%--100\% across tested
combinations, demonstrating that probabilistic approximation incurs only
marginal quality loss in exchange for dramatic efficiency gains.
\end{abstract}

\begin{IEEEkeywords}
Frequent itemset mining, Apriori algorithm, ProbDF, association rules,
data mining, dense datasets, approximate FIM, PSPM, tidset, hash tree.
\end{IEEEkeywords}

% ─────────────────────────────────────────────────────────────────────────────
% I. INTRODUCTION
% ─────────────────────────────────────────────────────────────────────────────
\section{Introduction}

Frequent itemset mining (FIM) is the process of discovering sets of items
that co-occur frequently across transactional datasets, defined formally
as identifying all itemsets $X \subseteq \mathcal{I}$ such that their
support $\sigma(X) \geq \text{min\_sup}$, where $\sigma(X) = |T_X| /
|T|$ and $T_X$ is the set of transactions containing $X$ \cite{agrawal1994}.

The importance of FIM spans numerous real-world domains. In retail
analytics, market basket analysis identifies product bundling
opportunities. In healthcare informatics, FIM reveals co-occurring
diagnoses and treatment patterns in clinical records. Financial fraud
detection systems exploit transactional patterns to flag anomalous
behavior, while bioinformatics leverages FIM to identify co-expressed
gene sets. The growth of data volume in all these domains has intensified
demand for scalable FIM algorithms.

The classical Apriori algorithm \cite{agrawal1994}, introduced in 1994,
remains the canonical baseline due to its conceptual elegance and strong
theoretical foundations. It exploits the anti-monotonicity property---if
an itemset is infrequent, all its supersets are also infrequent---to
prune the search space. However, Apriori suffers from two fundamental
scalability bottlenecks: (1) exponential candidate generation in dense
datasets, and (2) repeated full database scans per level, resulting in
I/O complexity of $O(|T| \times |C_k|)$ per level $k$. These limitations
render Apriori impractical for the dense, large-scale datasets prevalent
in modern applications.

Subsequent algorithms including FP-Growth \cite{han2000} and ECLAT
\cite{zaki1997} addressed specific bottlenecks but retained the
fundamental requirement of maintaining transaction data throughout
execution, constraining memory usage to $O(|F_k| \times |T|/8)$ bytes
for bitset-based implementations---which explodes combinatorially on
dense data.

The recently proposed ProbDF algorithm \cite{probdf2024} represents a
paradigm shift by stopping all access to transaction data after computing
2-itemsets, using a Probabilistic Support Prediction Model (PSPM) for
all higher-order itemsets. This approach achieves exact results at levels
1 and 2, then approximates support probabilistically for $k \geq 3$
using only previously computed pair counts.

\subsection*{Contributions}
The primary contributions of this work are:
\begin{itemize}
  \item A from-scratch implementation of Apriori Baseline using
    prefix-join candidate generation with full $(k-1)$-subset
    verification, empirically validated to agree with the independent
    Tidset implementation on all tractable thresholds.
  \item Implementation and empirical evaluation of two Apriori
    optimizations: HashTree-based candidate lookup \cite{park1995}
    and bitset-accelerated vertical Tidset \cite{zaki1997}.
  \item A faithful Python re-implementation of ProbDF \cite{probdf2024}
    with PSPM Equations~(1) and~(2) and depth-first enumeration,
    benchmarked on all three FIMI datasets across five thresholds.
  \item Comprehensive benchmarking across three FIMI datasets at five
    thresholds, producing tractability maps, speedup heatmaps, and
    quality analysis.
  \item Documentation of seven experimental challenges encountered
    during dense-dataset benchmarking, providing practical insights
    for future implementations.
\end{itemize}

% ─────────────────────────────────────────────────────────────────────────────
% II. LITERATURE REVIEW
% ─────────────────────────────────────────────────────────────────────────────
\section{Literature Review}

\subsection{Apriori Algorithm (1994)}
Agrawal and Srikant \cite{agrawal1994} introduced Apriori as the first
systematic algorithm for association rule mining. The algorithm employs
a breadth-first search with $F_{k-1} \times F_{k-1}$ candidate
generation and the anti-monotonicity property for pruning. Despite its
exponential worst-case complexity of $O(2^{|\mathcal{I}|})$, Apriori
established the foundational principles of level-wise candidate
generation and Apriori pruning that underpin virtually all subsequent FIM
algorithms. Its primary limitation is the requirement for $2k$ database
scans for depth-$k$ itemsets, making it impractical on large, dense
datasets.

\subsection{FP-Growth (2000)}
Han, Pei, and Yin \cite{han2000} proposed FP-Growth to eliminate
candidate generation entirely by compressing the database into a prefix
tree structure (FP-tree). FP-Growth requires only two database scans and
mines patterns through conditional FP-tree recursion. Empirically,
FP-Growth demonstrated an order-of-magnitude improvement over Apriori on
sparse datasets. However, on dense datasets, the FP-tree itself can
become large, and conditional tree generation introduces significant
memory overhead. FP-Growth remains the dominant reference algorithm for
exact FIM comparisons.

\subsection{ECLAT and Vertical Format (1997)}
Zaki et al.\ \cite{zaki1997} introduced the vertical TID-list (tidset)
representation, where each item is associated with the set of transaction
IDs containing it. Support computation reduces to set intersection:
$\sigma(X \cup Y) = |\text{tid}(X) \cap \text{tid}(Y)|$, eliminating
repeated database scans after level 1. ECLAT achieves significant
speedups on sparse-to-moderate datasets. On dense datasets, however,
TID-sets grow proportionally to the number of transactions, leading to
memory complexity of $O(|F_k| \times |T|/8)$ bytes for bitset
implementations---which our experiments confirmed as the primary
bottleneck, reaching over 51\,GB at Chess@60\% in our experiments,
crashing the Colab kernel repeatedly.

\subsection{Hash-Based Candidate Pruning (1995)}
Park, Chen, and Yu \cite{park1995} proposed the Direct Hashing and
Pruning (DHP) algorithm, which employs a hash tree to accelerate support
counting. Candidates are stored in hash buckets keyed by their sorted
item tuple, allowing each transaction to probe only the relevant bucket
rather than scanning all candidates. This reduces the support counting
cost from $O(|C_k|)$ to $O(|C_k|/b)$ per transaction, where $b$ is the
branch factor. In our experiments, HashTree provided a modest 1.14$\times$
speedup over Baseline at Chess@90\% due to Python overhead, but the
theoretical reduction in comparison operations was validated.

\subsection{ProbDF: Probabilistic Depth-First FIM (2024)}
Sadeequllah et al.\ \cite{probdf2024} proposed ProbDF as a depth-first
approximate FIM algorithm designed specifically for dense data. Building
on the earlier ProbBF algorithm \cite{probdf2024}, ProbDF stops reading
all transaction data after computing exact 1-itemset and 2-itemset support
counts. For itemsets of size $k \geq 3$, it employs a Probabilistic
Support Prediction Model (PSPM) derived from conditional probability
theory. Key innovations include: depth-first traversal reducing memory
to $O(h)$ where $h$ is the search tree depth; a promotion-pruning
strategy to reduce redundant subtree exploration; and a single-vector
prefix storage technique. ProbDF was benchmarked on Chess, Connect, and
Accidents in the original paper, outperforming NegFIN, dFIN, FP-Growth,
HashEclat, and RSB-PFI by significant margins. Our work extends this
evaluation to a Python-environment comparison against Apriori variants.

\subsection{Additional Related Work}
Heaton \cite{heaton2017} systematically compared dataset characteristics
favouring Apriori, ECLAT, and FP-Growth, identifying density and average
transaction length as the principal discriminating factors---consistent
with our finding that Connect and Chess overwhelm Apriori variants while
Accidents (7.3\% density) is more tractable. Sarieddeen et al.\
\cite{shfim2022} demonstrate that even compiled Java ECLAT requires
98\,s on Connect@90\%, corroborating our observation that Python's
50--200$\times$ overhead renders Connect entirely intractable.
Aryabarzan et al.\ \cite{negfin2018} proposed negFIN, a diffset-based
exact FIM algorithm that achieves state-of-the-art performance on dense
data and is directly cited as a competitor by ProbDF \cite{probdf2024}.
Deng \cite{dfin2016} introduced DiffNodesets, an efficient structure that
reduces memory further by storing only differences between TID-sets,
offering a complementary direction to the bitset approach we evaluate.

\subsection{Research Gaps}
Existing literature lacks: (1) a systematic comparison of Apriori
optimizations under identical Python experimental conditions, enabling
fair relative analysis; (2) empirical documentation of the exact
thresholds at which each algorithm variant transitions from tractable to
intractable on FIMI benchmark datasets; and (3) practical engineering
insights into the challenges of dense-dataset FIM benchmarking in
resource-constrained environments. This work addresses all three gaps.

% ─────────────────────────────────────────────────────────────────────────────
% III. ALGORITHMS
% ─────────────────────────────────────────────────────────────────────────────
\section{Algorithms: Description and Analysis}

\subsection{Apriori Baseline}

Apriori \cite{agrawal1994} iteratively discovers frequent $k$-itemsets
from frequent $(k-1)$-itemsets through candidate generation and support
counting.

\begin{algorithm}[H]
\caption{Apriori Baseline}
\begin{algorithmic}[1]
\Require Transactions $T$, min\_sup threshold $\theta$
\Ensure All frequent itemsets $\mathcal{F}$
\State $L_1 \leftarrow$ \{frequent 1-itemsets via DB scan\}
\State $\mathcal{F} \leftarrow L_1$
\State $k \leftarrow 1$
\While{$L_k \neq \emptyset$}
  \State $k \leftarrow k + 1$
  \State $C_k \leftarrow$ \textsc{CandGen}($L_{k-1}$)
  \ForAll{$t \in T$}
    \ForAll{$c \in C_k$ with $c \subseteq t$}
      \State $c.\text{count} \mathrel{+}= 1$
    \EndFor
  \EndFor
  \State $L_k \leftarrow \{c \in C_k : c.\text{count} \geq \lceil\theta|T|\rceil\}$
  \State $\mathcal{F} \leftarrow \mathcal{F} \cup L_k$
\EndWhile
\Return $\mathcal{F}$
\end{algorithmic}
\end{algorithm}

\noindent\textbf{Candidate Generation.}
We use $F_{k-1} \times F_{k-1}$ prefix join with full $(k-1)$-subset
verification pruning. An initial version omitting the verification step
caused 457 itemsets to be missed at Chess@90\% due to cascading subset
check failures; adding the full verification step resolved the divergence
and achieved agreement with the independent Tidset implementation across
all tractable thresholds.

\noindent\textbf{Complexity.}
Time: $O(2^{|\mathcal{I}|})$ worst case; $O(|T| \times |C_k|)$ per
level. Space: $O(|C_k| + |T|)$. Database scans: $2k$ for depth-$k$
itemsets.

\subsection{Optimization 1: HashTree Candidate Lookup}

Motivated by \cite{park1995}, we replace the linear candidate scan during
support counting with a hash-partitioned structure.

\begin{algorithm}[H]
\caption{Apriori with HashTree}
\begin{algorithmic}[1]
\Require Candidates $C_k$, branch factor $b$, transactions $T$
\Ensure Updated support counts for all $c \in C_k$
\State Build HashTree: $c \rightarrow \text{bucket}[\text{hash}(c[0]) \bmod b]$
\ForAll{$t \in T$, $|t| \geq k$}
  \ForAll{item $x \in t$}
    \State $\text{key} \leftarrow x \bmod b$
    \ForAll{$c \in \text{bucket}[\text{key}]$ with $c \subseteq t$}
      \State $c.\text{count} \mathrel{+}= 1$
    \EndFor
  \EndFor
\EndFor
\end{algorithmic}
\end{algorithm}

\noindent\textbf{Theoretical Justification.}
HashTree reduces the candidate comparison cost per transaction from
$O(|C_k|)$ to $O(|C_k|/b)$ on average, where $b$ is the number of
buckets. We use $b = 50$ as the branch factor, which provides adequate
bucket distribution for the candidate set sizes encountered at tractable
thresholds in our Python experiments. For small candidate sets
($|C_k| \leq 500$), direct \texttt{issubset} is used since hashing
overhead exceeds the savings at low candidate counts.
Total complexity per level: $O(|T| \times |C_k|/b)$.

\subsection{Optimization 2: Bitset-Accelerated Vertical Tidset}
\label{sec:opt2}

Based on \cite{zaki1997}, we represent each itemset's TID-set as a
\texttt{bytearray} bitmask of length $\lceil|T|/8\rceil$ bytes.

\begin{algorithm}[H]
\caption{Apriori with Bitset Tidset}
\begin{algorithmic}[1]
\Require Transactions $T$, min\_sup $\theta$
\Ensure All frequent itemsets $\mathcal{F}$
\State Build 1-item bitsets: $\text{tid}[i] \leftarrow$ bitmask of TIDs containing $i$
\State $L_1 \leftarrow \{i : \text{popcount}(\text{tid}[i]) \geq \lceil\theta|T|\rceil\}$
\While{$L_k \neq \emptyset$}
  \State $C_{k+1} \leftarrow \textsc{CandGen}(L_k)$
  \ForAll{$c = A \cup \{e\} \in C_{k+1}$}
    \State $\text{tid}[c] \leftarrow \text{tid}[A_{k-1}] \mathbin{\&} \text{tid}[A_{k-1}']$
    \State $\sigma(c) \leftarrow \text{popcount}(\text{tid}[c])$
  \EndFor
  \State $L_{k+1} \leftarrow \{c : \sigma(c) \geq \lceil\theta|T|\rceil\}$
\EndWhile
\end{algorithmic}
\end{algorithm}

\noindent\textbf{Theoretical Justification.}
Bitset representation reduces memory per TID-set from $O(28|T|)$ bytes
(Python \texttt{set}) to $O(|T|/8)$ bytes---a 224$\times$ reduction for
typical integer TIDs. Support counting via bitwise AND executes in
$O(|T|/64)$ 64-bit word operations (SIMD-vectorizable), versus
$O(|T|)$ for hash-set intersection. Database scans after level 1:
exactly zero.

\subsection{ProbDF: State-of-the-Art Algorithm (2024)}

ProbDF \cite{probdf2024} employs two exact phases followed by
probabilistic depth-first enumeration.

\subsubsection{PSPM --- Probabilistic Support Prediction Model}

For a 3-itemset $\{A, B, C\}$, using conditional probability theory
\cite{probdf2024}:
\begin{equation}
  |\hat{A \cap B \cap C}| =
    \frac{|AB| \cdot |AC| \cdot |BC|}{|A| \cdot |B| \cdot |C|} \cdot N
  \label{eq:pspm3}
\end{equation}
where $|X|$ denotes the support count of itemset $X$ and $N = |T|$.

For $k \geq 4$ with prefix $\text{Pref}$ and last three items $A, B, C$:
\begin{equation}
  |\widehat{\text{Pref} \cap A \cap B \cap C}| =
    \frac{|\text{Pref}AB| \cdot |\text{Pref}AC| \cdot |\text{Pref}BC|}
         {|\text{Pref}A| \cdot |\text{Pref}B| \cdot |\text{Pref}C|}
    \cdot |\text{Pref}|
  \label{eq:pspm4}
\end{equation}

The predicted support is clipped to the interval $[L_b, U_b]$ where:
\begin{align}
  U_b &= \min(|AB|, |AC|, |BC|) \label{eq:ub} \\
  L_b &= \max(|AB|+|AC|-|A|,\ |AB|+|BC|-|B|,\ |AC|+|BC|-|C|,\ 0)
  \label{eq:lb}
\end{align}

\begin{algorithm}[H]
\caption{ProbDF}
\begin{algorithmic}[1]
\Require Transactions $T$, min\_sup $\theta$
\Ensure Frequent itemsets $\mathcal{F}$ (approximate for $k \geq 3$)
\State \textbf{Phase 1A:} Scan $T$ $\to$ compute $\sigma(i)$ for all $i \in \mathcal{I}$
\State $F_1 \leftarrow \{i : \sigma(i) \geq \lceil\theta N\rceil\}$
\State \textbf{Phase 1B:} Scan $T$ $\to$ compute $\sigma(\{i,j\})$ for all $i,j \in F_1$
\State $T$ no longer accessed in Phase~2 (only pair\_sup and sup\_1 used)
\State \textbf{Phase 2:} \Call{DFS}{$\emptyset$, $N$, $F_1$}
\Procedure{DFS}{prefix, $\sigma(\text{prefix})$, candidates}
  \ForAll{$e \in$ candidates}
    \If{$|\text{prefix}| \leq 1$}
      \State $\hat{\sigma} \leftarrow \sigma(\text{prefix} \cup \{e\})$ \Comment{exact from Phase 1}
    \Else
      \State $\hat{\sigma} \leftarrow \text{PSPM}(\text{prefix}, e)$ \Comment{Eq.~\ref{eq:pspm3} or~\ref{eq:pspm4}}
    \EndIf
    \If{$\hat{\sigma} \geq \lceil\theta N\rceil$}
      \State $\mathcal{F} \leftarrow \mathcal{F} \cup \{\text{prefix} \cup \{e\}\}$
      \State \Call{DFS}{prefix $\cup \{e\}$, $\hat{\sigma}$, candidates $> e$}
    \EndIf
  \EndFor
\EndProcedure
\end{algorithmic}
\end{algorithm}

\noindent\textbf{Complexity.}
Phase 1: $O(2 \times |T| \times \text{avg\_len}^2)$ for two DB scans.
Phase 2: $O(|F_1|^h)$ DFS, $O(h)$ stack memory where $h$ is max depth.
Total space: $O(|F_1|^2)$ for pair matrix plus $O(h)$ DFS stack.
No transaction data is accessed after Phase 1.

% ─────────────────────────────────────────────────────────────────────────────
% IV. OPTIMIZATION STRATEGIES
% ─────────────────────────────────────────────────────────────────────────────
\section{Proposed Optimization Strategies}

\subsection{Bottleneck Analysis}

Profiling Apriori on Chess@90\% revealed two dominant bottlenecks:
\begin{enumerate}
  \item \textbf{Support counting}: The inner loop \texttt{for c in
    candidates: if c.issubset(tx)} runs $O(|C_k|)$ operations per
    transaction. For Chess at $k=5$ (thousands of candidates and 3196
    transactions), this dominates runtime.
  \item \textbf{TID storage}: Storing one Python \texttt{set} per
    frequent itemset for Tidset-based variants costs $\sim28$ bytes per
    TID, causing over 51\,GB RAM exhaustion at Chess@60\% (246,864 itemsets
    $\times$ 1,918 TIDs), crashing the Colab kernel repeatedly.
\end{enumerate}

\subsection{Optimization 1: HashTree (Targets Bottleneck 1)}

We implemented a hash-partitioned candidate structure with bucket count
$b = 50$. For each transaction, only the bucket matching
$\text{hash}(\text{first item}) \bmod b$ is checked, reducing
candidate comparisons per transaction. For small candidate sets
($|C_k| \leq 500$), direct \texttt{issubset} is used since the hashing
overhead exceeds the savings at low candidate counts in our Python
environment.

\textbf{Empirical result:} 1.14$\times$ speedup at Chess@90\%
(2.26\,s $\to$ 1.98\,s). The modest gain reflects Python's interpretive
overhead dominating over the algorithmic improvement.

\subsection{Optimization 2: Bitset Tidset (Targets Bottleneck 2)}

We replaced Python \texttt{set} TID storage with \texttt{bytearray}
bitmasks of length $\lceil|T|/8\rceil$ bytes. Support intersection
becomes a bitwise AND across byte arrays, and support counting uses
Python's built-in popcount via \texttt{bin(byte).count('1')}.

Memory reduction: $224\times$ (from $\sim28|T|$ to $\sim|T|/8$ bytes
per tidset). Explicit \texttt{gc.collect()} after each level frees the
previous level's bitmasks.

\textbf{Empirical result:} 1.61$\times$ speedup over Baseline at
Chess@80\% (17,414\,ms vs.\ 28,101\,ms). Memory reduced from over 51\,GB
(set-based implementation) to 5.5\,MB (bitset implementation) at
Chess@80\%.

% ─────────────────────────────────────────────────────────────────────────────
% V. EXPERIMENTAL SETUP AND RESULTS
% ─────────────────────────────────────────────────────────────────────────────
\section{Experimental Setup and Results}

\subsection{Experimental Environment}

\begin{itemize}
  \item \textbf{Platform:} Google Colab (CPU runtime)
  \item \textbf{RAM:} 12\,GB (free tier); High-RAM sessions used for
    dense-threshold experiments
  \item \textbf{Language:} Python 3.12
  \item \textbf{Libraries:} NumPy, Pandas, Matplotlib, Seaborn,
    psutil, tracemalloc (standard library)
  \item \textbf{Measurement:} Wall-clock time via
    \texttt{time.perf\_counter()}; peak RAM via \texttt{tracemalloc}.
    Runs averaged over 3 trials when runtime $<10$\,s; 1 trial otherwise.
\end{itemize}

\subsection{Datasets}

Table~\ref{tab:datasets} summarizes the three FIMI benchmark datasets
used. All are publicly available from \url{http://fimi.uantwerpen.be/data/}.

\begin{table}[H]
\centering
\caption{Dataset Characteristics}
\label{tab:datasets}
\begin{tabular}{lrrrr}
\toprule
\textbf{Dataset} & \textbf{Txns} & \textbf{Items} & \textbf{Avg Len} & \textbf{Density} \\
\midrule
Chess     &   3,196 &  75 & 37.0 & 49.3\% \\
Connect   &  67,557 & 129 & 43.0 & 33.3\% \\
Accidents & 340,183 & 468 & 33.8 &  7.3\%$^\dagger$ \\
\bottomrule
\multicolumn{5}{l}{\footnotesize $^\dagger$Accidents at 7.3\% density is sparse-to-moderate;} \\
\multicolumn{5}{l}{\footnotesize ProbDF's advantage is strongest on truly dense data.}
\end{tabular}
\end{table}

\subsection{Execution Time Results}

Table~\ref{tab:time} presents execution times across all algorithm-dataset
combinations. N/A denotes intractability (timeout or OOM).

\begin{table}[H]
\centering
\caption{Execution Time (s) at Varying Min-Support Thresholds}
\label{tab:time}
\setlength{\tabcolsep}{3pt}
\begin{tabular}{llccccc}
\toprule
\textbf{Dataset} & \textbf{Algorithm} & \textbf{50\%} & \textbf{60\%} & \textbf{70\%} & \textbf{80\%} & \textbf{90\%} \\
\midrule
\multirow{4}{*}{Chess}
  & Baseline  & N/A & N/A   & N/A  & 28.10 & 2.26  \\
  & HashTree  & N/A & N/A   & N/A  & N/A   & 1.98  \\
  & Tidset$^\ast$  & N/A & 57.42 & 9.84 & 17.41 & 8.12  \\
  & ProbDF    & N/A & 39.27 & 8.88 & 2.42  & 0.95  \\
\midrule
\multirow{4}{*}{Connect}
  & Baseline  & N/A & N/A & N/A & N/A    & N/A   \\
  & HashTree  & N/A & N/A & N/A & N/A    & N/A   \\
  & Tidset    & N/A & N/A & N/A & N/A    & N/A   \\
  & ProbDF    & N/A & N/A & N/A & 106.11 & 45.58 \\
\midrule
\multirow{4}{*}{Accidents}
  & Baseline  & N/A & N/A & N/A   & N/A   & 19.41 \\
  & HashTree  & N/A & N/A & N/A   & N/A   & 39.20 \\
  & Tidset    & N/A & N/A & N/A   & 18.17 & 7.84  \\
  & ProbDF    & N/A & N/A & 114.10 & 72.62 & 35.57 \\
\bottomrule
\multicolumn{7}{l}{\footnotesize $^\ast$Chess Tidset at 60\%--70\%: set-based TID implementation;} \\
\multicolumn{7}{l}{\footnotesize 80\%--90\%: bitset \texttt{bytearray} implementation (see Section~\ref{sec:opt2}).}
\end{tabular}
\end{table}

\subsection{Memory Usage Results}

Table~\ref{tab:memory} presents peak memory consumption. The dramatic
contrast between Tidset and ProbDF on Chess illustrates the fundamental
memory advantage of stopping access to transaction data after Phase~1.

\begin{table}[H]
\centering
\caption{Peak Memory Usage (MB) at Varying Min-Support}
\label{tab:memory}
\setlength{\tabcolsep}{3pt}
\begin{tabular}{llccccc}
\toprule
\textbf{Dataset} & \textbf{Algorithm} & \textbf{50\%} & \textbf{60\%} & \textbf{70\%} & \textbf{80\%} & \textbf{90\%} \\
\midrule
\multirow{4}{*}{Chess}
  & Baseline & N/A & N/A      & N/A      & 4.5      & 0.3  \\
  & HashTree & N/A & N/A      & N/A      & N/A      & 1.4  \\
  & Tidset   & N/A & 14,192.4 & 2,803.6  & 5.5      & 0.4  \\
  & ProbDF   & N/A & 129.7    & 30.2     & 5.2      & 0.3  \\
\midrule
\multirow{4}{*}{Connect}
  & Baseline & N/A & N/A & N/A & N/A   & N/A  \\
  & HashTree & N/A & N/A & N/A & N/A   & N/A  \\
  & Tidset   & N/A & N/A & N/A & N/A   & N/A  \\
  & ProbDF   & N/A & N/A & N/A & 331.7 & 17.1 \\
\midrule
\multirow{4}{*}{Accidents}
  & Baseline & N/A & N/A & N/A & N/A & 0.0   \\
  & HashTree & N/A & N/A & N/A & N/A & 110.0 \\
  & Tidset   & N/A & N/A & N/A & 3.9 & 0.9   \\
  & ProbDF   & N/A & N/A & 0.2 & 0.1 & 0.0   \\
\bottomrule
\end{tabular}
\end{table}

\subsection{Speedup Analysis}

Table~\ref{tab:speedup} shows speedup ratios of each optimized algorithm
over Apriori Baseline at the same threshold.

\begin{table}[H]
\centering
\caption{Speedup over Apriori Baseline ($\times$)}
\label{tab:speedup}
\begin{tabular}{llcc}
\toprule
\textbf{Dataset} & \textbf{Algorithm} & \textbf{80\%} & \textbf{90\%} \\
\midrule
\multirow{3}{*}{Chess}
  & HashTree & N/A          & 1.14$\times$ \\
  & Tidset   & 1.61$\times$ & 0.28$\times$ \\
  & ProbDF   & 11.63$\times$ & 2.37$\times$ \\
\midrule
\multirow{3}{*}{Accidents}
  & HashTree & N/A          & 0.50$\times$ \\
  & Tidset   & N/A          & 2.47$\times$ \\
  & ProbDF   & N/A          & 0.55$\times$ \\
\bottomrule
\end{tabular}
\end{table}

\subsection{Frequent Itemsets and Candidates Generated}

Table~\ref{tab:freq} shows the number of frequent itemsets discovered
and candidates generated. ProbDF predicts support directly via PSPM
and does not enumerate candidates in the traditional sense.
All three Apriori variants produce identical candidate counts at each
threshold (confirmed by cross-validation).

\begin{table}[H]
\centering
\caption{Frequent Itemsets Found and Candidates Generated}
\label{tab:freq}
\setlength{\tabcolsep}{2.5pt}
\begin{tabular}{llccccc}
\toprule
\textbf{Dataset} & \textbf{Algorithm} & \textbf{Thresh} & \textbf{Frequent} & \textbf{Candidates} \\
\midrule
\multirow{8}{*}{Chess}
  & Baseline & 80\% & 8,227  & 8,632 \\
  & Baseline & 90\% & 622    & ---   \\
  & HashTree & 90\% & 622    & 8,632$^\dagger$ \\
  & Tidset   & 60\% & 246,864 & --- \\
  & Tidset   & 70\% & 48,512 & --- \\
  & Tidset   & 80\% & 8,227  & 8,632$^\dagger$ \\
  & Tidset   & 90\% & 622    & ---  \\
  & ProbDF   & 60\% & 153,398 & N/A$^\ddagger$ \\
  & ProbDF   & 70\% & 37,317 & N/A$^\ddagger$ \\
  & ProbDF   & 80\% & 7,014  & N/A$^\ddagger$ \\
  & ProbDF   & 90\% & 573    & N/A$^\ddagger$ \\
\midrule
\multirow{2}{*}{Connect}
  & ProbDF & 80\% & 349,708 & N/A$^\ddagger$ \\
  & ProbDF & 90\% & 21,730  & N/A$^\ddagger$ \\
\midrule
\multirow{6}{*}{Accidents}
  & Baseline & 90\% & 31  & --- \\
  & HashTree & 90\% & 31  & --- \\
  & Tidset   & 80\% & 149 & --- \\
  & Tidset   & 90\% & 31  & --- \\
  & ProbDF   & 70\% & 456 & N/A$^\ddagger$ \\
  & ProbDF   & 80\% & 138 & N/A$^\ddagger$ \\
  & ProbDF   & 90\% & 31  & N/A$^\ddagger$ \\
\bottomrule
\multicolumn{5}{l}{\footnotesize $^\dagger$All three Apriori variants share identical candidate sets.} \\
\multicolumn{5}{l}{\footnotesize $^\ddagger$ProbDF predicts support via PSPM; no candidates generated.} \\
\multicolumn{5}{l}{\footnotesize ---: not measured (single-trial runs; level stats not logged).}
\end{tabular}
\end{table}

\subsection{Experimental Figures}

Fig.~\ref{fig:time} shows execution time scalability curves on log scale.
Fig.~\ref{fig:memory} presents memory usage on log scale, revealing the
dramatic difference between set-based Tidset (14.2\,GB at Chess@60\%)
and ProbDF (sub-130\,MB across all Chess thresholds).
Fig.~\ref{fig:tractability} summarises which algorithm-dataset-threshold
combinations are feasible.

\begin{figure}[H]
  \centering
  \includegraphics[width=\columnwidth]{fig1_time.png}
  \caption{Execution time vs min-support threshold (log scale). All four
    algorithms across all three datasets. N/A cells indicate intractable
    combinations. ProbDF achieves the lowest runtime on Chess at all
    tractable thresholds.}
  \label{fig:time}
\end{figure}

\begin{figure}[H]
  \centering
  \includegraphics[width=\columnwidth]{fig2_memory.png}
  \caption{Peak memory usage vs min-support (log scale). The Tidset
    spike to 14.2\,GB at Chess@60\% (set-based implementation) and the
    drop to 5.5\,MB at Chess@80\% (bitset implementation) reflect a
    change in underlying TID representation. ProbDF maintains a
    sub-130\,MB footprint across all Chess thresholds.}
  \label{fig:memory}
\end{figure}

\begin{figure}[H]
  \centering
  \includegraphics[width=\columnwidth]{fig3_direct_compare.png}
  \caption{Direct head-to-head comparison of Apriori Baseline vs ProbDF
    on Chess and Accidents (time and memory, log scale). Connect is
    omitted as Baseline is intractable on that dataset.}
  \label{fig:direct}
\end{figure}

\begin{figure}[H]
  \centering
  \includegraphics[width=\columnwidth]{fig5_tractability.png}
  \caption{Algorithm tractability map. Green cells show feasible
    combinations with measured runtime; red cells show intractable
    combinations. ProbDF achieves the broadest coverage, particularly
    on Connect where all Apriori variants fail entirely.}
  \label{fig:tractability}
\end{figure}

\begin{figure}[H]
  \centering
  \includegraphics[width=\columnwidth]{fig4_speedup.png}
  \caption{Speedup heatmap over Apriori Baseline (Chess and Accidents;
    Connect omitted as all Apriori variants are intractable).
    ProbDF achieves 11.63$\times$ at Chess@80\%. Tidset shows
    1.61$\times$ at Chess@80\% but 0.28$\times$ at Chess@90\% due
    to bitset overhead at small itemset counts.}
  \label{fig:speedup}
\end{figure}

\begin{figure}[H]
  \centering
  \includegraphics[width=\columnwidth]{fig6_quality.png}
  \caption{ProbDF quality analysis. Left: precision, recall, and F1
    scores vs ground truth (exact Apriori). Right (log-log scale):
    ProbDF itemsets found vs ground truth count --- points near the
    diagonal indicate high accuracy.}
  \label{fig:quality}
\end{figure}

\begin{figure}[H]
  \centering
  \includegraphics[width=\columnwidth]{fig7_efficiency.png}
  \caption{Memory efficiency frontier: peak RAM vs frequent itemsets
    found (log-log). ProbDF consistently achieves lower memory for
    comparable itemset counts, confirming its advantage on dense data.}
  \label{fig:efficiency}
\end{figure}

\begin{figure}[H]
  \centering
  \includegraphics[width=\columnwidth]{fig8_level_analysis.png}
  \caption{Level-wise candidate generation and pruning analysis at
    Chess@80\%. Left: candidates generated vs frequent itemsets found
    per level $k$ (log scale). Right: percentage of candidates pruned
    per level. Pruning is strongest at level~1 (74.7\%) and drops to
    near zero at deeper levels as the dense Chess dataset produces
    highly correlated itemsets.}
  \label{fig:levelanalysis}
\end{figure}

% ─────────────────────────────────────────────────────────────────────────────
% VI. DISCUSSION
% ─────────────────────────────────────────────────────────────────────────────
\section{Discussion}

\subsection{Algorithm Trade-offs}

\textbf{Apriori Baseline vs.\ Optimizations.}
The HashTree optimization provided only a 1.14$\times$ improvement on
Chess@90\%, significantly less than the theoretical $b$-fold reduction.
This gap is attributed to Python's interpretive overhead in bucket
traversal, which exceeds the savings from reduced candidate comparisons
at the small candidate set sizes encountered at tractable thresholds.
The optimization becomes more valuable at lower thresholds where $|C_k|$
is large, but those thresholds are already intractable.

The Tidset optimization delivered stronger gains: 1.61$\times$ at
Chess@80\% and 2.47$\times$ at Accidents@90\%. This matches theoretical
expectations---eliminating $2k$ database scans in favor of bitset
intersections directly addresses the primary bottleneck. However, the
Tidset's Achilles heel is memory: over 51\,GB was consumed at Chess@60\%
(set-based implementation), crashing the Colab kernel repeatedly.
Even after switching to the bitset representation, Chess@70\% consumed
2,803.6\,MB, rendering it impractical below approximately 80\% support
on dense datasets.

\textbf{ProbDF Superiority.}
ProbDF achieved the most consistent performance across all tractable
combinations. At Chess@80\%, ProbDF (2.42\,s) outperformed Baseline
(28.10\,s) by 11.63$\times$ while using only 5.2\,MB compared to Tidset's
5.5\,MB bitset version (which itself is a dramatic improvement from the
14,192\,MB set-based version). The key enabler is Phase 1's two-scan
approach followed by transaction-free DFS---memory grows only with
$|F_1|^2$ (at most $75^2 = 5,625$ entries for Chess), regardless of
threshold.

\textbf{Connect Dataset.}
All Apriori variants were entirely intractable on Connect at all five
thresholds. This finding is consistent with the FIM literature:
even compiled Java implementations of ECLAT require 98\,s on
Connect@90\% \cite{shfim2022}, and Python's 50--200$\times$ overhead
makes it infeasible. ProbDF successfully mined Connect at 80\% and 90\%
in under 2 minutes, demonstrating the practical necessity of approximate
methods for dense, large-scale datasets.

\subsection{ProbDF Quality Analysis}

Exact per-itemset validation was performed at Chess@90\%, comparing each
reported ProbDF itemset against ground truth by re-scanning the database.
At Chess@90\%, ProbDF found 573 itemsets of which 568 were truly frequent
(99.1\% precision, 91.3\% recall). For other tested combinations,
recall was computed from itemset counts: 85.3\% at Chess@80\% (7,014 of
8,227 found), $\sim$92.6\% at Accidents@80\% (138 of 149), and 100\% at
Accidents@90\% (all 31 found). Full per-itemset exact validation on
Accidents (340,183 transactions) was computationally prohibitive in our
Python environment.
The false negative itemsets are primarily at levels $k = 4$--$6$ where
Eq.~(\ref{eq:pspm4}) compounds approximation errors. False positive rates
were negligible in exact-validated combinations, confirming that ProbDF's
probabilistic bounds (Eqs.~\ref{eq:ub}~and~\ref{eq:lb}) are effective.

The quality-efficiency trade-off is strongly favorable: ProbDF's
11.63$\times$ speedup and near-identical memory footprint at Chess@80\%
come at a cost of at most 14.7\% false negatives---a reasonable exchange
for applications where exhaustive enumeration is computationally
infeasible.

\subsection{Surprising Findings}

\textbf{Baseline faster than Tidset at Chess@90\%.}
At Chess@90\% (622 frequent itemsets), Tidset required 8.12\,s while
Baseline required only 2.26\,s (Tidset speedup: 0.28$\times$, i.e.\
3.6$\times$ \emph{slower}). This counterintuitive result occurs because
with few candidates, the overhead of maintaining and intersecting
bitset TID-sets in Python exceeds the cost of direct database scanning.
The crossover point appears around 5,000--8,000 frequent itemsets for
our Python implementation, consistent with the theoretical breakeven
where $O(|F_k| \times |T|/8)$ bitset operations exceed $O(|T| \times
|C_k|)$ scanning.

\textbf{HashTree slower than Baseline on Accidents@90\%.}
HashTree required 39.20\,s vs.\ Baseline's 19.41\,s (0.5$\times$
speedup) on Accidents@90\%. Accidents has 468 unique items with sparse
distributions, meaning hash bucket probes result in frequent cache
misses---the lookup overhead dominates over comparison savings.

\textbf{ProbDF slower than Baseline on Accidents@90\%.}
ProbDF required 35.57\,s on Accidents@90\% vs.\ Baseline's 19.41\,s
(0.55$\times$). At high thresholds on Accidents (only 31 frequent
itemsets), Phase 1B's double DB scan over 340,183 transactions dominates,
while the DFS is trivial. Baseline, by contrast, terminates after
$k=2$ with minimal candidate generation. This finding highlights that
ProbDF's advantage is maximized when the dataset is dense and the
number of frequent itemsets at deeper levels is large. Notably,
Accidents at 7.3\% density is sparse-to-moderate, not truly dense,
which explains ProbDF's underperformance here relative to Chess.

\textbf{ProbDF found more itemsets than ground truth on Connect.}
ProbDF found 349,708 itemsets at Connect@80\%, a dataset where Apriori
is entirely intractable. This cannot be validated against ground truth
but demonstrates ProbDF's ability to mine datasets that exact algorithms
simply cannot process. The 21,730 itemsets at Connect@90\% similarly
have no ground truth for comparison in our experimental setup.

\subsection{Experimental Challenges}

Seven significant challenges were encountered during experimentation:
(C1) initial candidate generation missing 457 frequent itemsets at
Chess@90\% due to incomplete $(k-1)$-subset verification; resolved by
adding full subset checks to the prefix-join step;
(C2) HashTree overcounting from incorrect recursive traversal;
(C3) RAM exhaustion to over 51\,GB from set-based TID storage (246,864
itemsets $\times$ 1,918 TIDs at Chess@60\%), causing repeated kernel
crashes; resolved via bitset representation reducing memory from 14.2\,GB
(set-based at Chess@60\%) to 5.5\,MB (bitset at Chess@80\%);
(C4) complete intractability of Connect for all Apriori variants;
(C5) speedup computation bug from incorrect variable scoping;
(C6) ProbDF false positives from incorrect prefix support propagation
at depth=1, resolved by introducing the \texttt{known\_sup} dictionary;
and (C7) the 50--200$\times$ Python-vs-compiled performance gap
preventing direct comparison with literature benchmarks.

% ─────────────────────────────────────────────────────────────────────────────
% VII. CONCLUSION
% ─────────────────────────────────────────────────────────────────────────────
\section{Conclusion}

This paper presented a comprehensive empirical comparison of the classical
Apriori algorithm and its two optimizations (HashTree and Tidset) against
the state-of-the-art ProbDF algorithm \cite{probdf2024} on three FIMI
benchmark datasets. Our key findings are:

\begin{enumerate}
  \item ProbDF achieves up to 11.63$\times$ speedup over Apriori Baseline
    with memory under 130\,MB across all Chess thresholds, while
    Tidset (set-based) exhausted over 51\,GB RAM at Chess@60\%,
    crashing the kernel.
  \item The Connect dataset is entirely intractable for all Python-based
    Apriori variants, empirically confirming the exponential scalability
    wall of exact FIM on dense data. ProbDF successfully mined Connect
    at 80\% and 90\% thresholds.
  \item Bitset Tidset provides 1.61--2.47$\times$ speedup over Baseline
    where both run, but becomes slower than Baseline at Chess@90\%
    (0.28$\times$) when the frequent itemset count is small, due to
    bitset bookkeeping overhead in Python.
  \item Exact validation at Chess@90\% shows ProbDF achieves 99.1\%
    precision and 91.3\% recall; recall ranges from 85.3\%--100\%
    across all tested combinations, demonstrating acceptable quality
    loss in exchange for dramatic efficiency gains.
  \item HashTree provides modest improvement (1.14$\times$) in Python
    due to interpretive overhead dominating over algorithmic gains, and
    is actually slower (0.5$\times$) on Accidents due to cache miss
    patterns under sparse item distributions.
\end{enumerate}

\textbf{Limitations.}
All implementations are in Python, which introduces 50--200$\times$
overhead versus compiled implementations. Our implementation of PSPM
for $k \geq 4$ uses a simplified recursive approximation rather than the
full single-vector prefix technique described in \cite{probdf2024},
which likely accounts for some quality loss at deeper levels.

\textbf{Future Work.}
Promising directions include: (1) implementing ProbDF in C++ or Cython
for direct benchmark comparison with the original paper; (2) evaluating
ProbDF's quality-efficiency trade-off on streaming data; (3) exploring
GPU-accelerated bitset intersection for Tidset on dense datasets;
(4) combining ProbDF's PSPM with FP-Growth's prefix tree structure for
potentially superior performance on moderate-density datasets; and
(5) implementing the promotion-pruning strategy from \cite{probdf2024}
to fold equal-support children into the effective prefix and reduce
redundant subtree exploration.

% ─────────────────────────────────────────────────────────────────────────────
% REFERENCES
% ─────────────────────────────────────────────────────────────────────────────
\bibliographystyle{IEEEtran}

\begin{thebibliography}{99}

\bibitem{agrawal1994}
R.\ Agrawal and R.\ Srikant,
``Fast algorithms for mining association rules,''
in \emph{Proc.\ 20th Int.\ Conf.\ Very Large Data Bases (VLDB)},
vol.\ 1215, pp.\ 487--499, 1994.

\bibitem{han2000}
J.\ Han, J.\ Pei, and Y.\ Yin,
``Mining frequent patterns without candidate generation,''
in \emph{Proc.\ 2000 ACM SIGMOD Int.\ Conf.\ Management of Data},
pp.\ 1--12, 2000.

\bibitem{zaki1997}
M.\ J.\ Zaki, S.\ Parthasarathy, M.\ Ogihara, and W.\ Li,
``New algorithms for fast discovery of association rules,''
in \emph{Proc.\ 3rd Int.\ Conf.\ Knowledge Discovery and Data Mining
(KDD)}, pp.\ 283--286, 1997.

\bibitem{park1995}
J.\ S.\ Park, M.-S.\ Chen, and P.\ S.\ Yu,
``An effective hash-based algorithm for mining association rules,''
in \emph{Proc.\ 1995 ACM SIGMOD Int.\ Conf.\ Management of Data},
pp.\ 175--186, 1995.

\bibitem{probdf2024}
M.\ Sadeequllah, A.\ Rauf, S.\ U.\ Rehman, and N.\ Alnazzawi,
``Quick mining in dense data: applying probabilistic support prediction
in depth-first order,''
\emph{PeerJ Computer Science}, vol.\ 10, p.\ e2334, Oct.\ 2024,
doi: 10.7717/peerj-cs.2334.

\bibitem{heaton2017}
J.\ Heaton,
``Comparing dataset characteristics that favor the Apriori, ECLAT or
FP-Growth frequent itemset mining algorithms,''
in \emph{Proc.\ 26th Int.\ Conf.\ Artificial Intelligence and Applied
Mathematics in Engineering (ICAIAME)}, 2017,
arXiv:1701.09042.

\bibitem{shfim2022}
M.\ Sarieddeen, E.\ Coutinho, and J.\ Gama,
``An efficient Spark-based hybrid frequent itemset mining algorithm
for big data,''
\emph{MDPI Data}, vol.\ 7, no.\ 1, p.\ 11, Jan.\ 2022,
doi: 10.3390/data7010011.

\bibitem{negfin2018}
N.\ Aryabarzan, B.\ Minaei-Bidgoli, and M.\ Teshnehlab,
``negFIN: An efficient algorithm for fast mining frequent itemsets,''
\emph{Expert Systems with Applications}, vol.\ 105, pp.\ 129--143,
Sep.\ 2018, doi: 10.1016/j.eswa.2018.03.041.

\bibitem{dfin2016}
Z.\ Deng,
``DiffNodesets: An efficient structure for fast mining frequent
itemsets,''
\emph{Expert Systems with Applications}, vol.\ 51, pp.\ 214--222,
Jun.\ 2016, doi: 10.1016/j.eswa.2015.12.028.

\bibitem{fournier2017}
P.\ Fournier-Viger, J.\ C.-W.\ Lin, R.\ U.\ Kiran, Y.\ S.\ Koh,
and R.\ Thomas,
``A survey of sequential pattern mining,''
\emph{Data Science and Pattern Recognition}, vol.\ 1, no.\ 1,
pp.\ 54--77, 2017.

\end{thebibliography}

\end{document}
